% SEGMENT OLP-0631-S01
% Part: history
% Chapter: biographies
% Section: alan-turing

\documentclass[../../../include/open-logic-section]{subfiles}

\begin{document}

\olfileid{his}{bio}{tur}

\olsection{ஆலன் டியூரிங்}

ஆலன் டியூரிங் 1912 ஜூன் 23 அன்று லண்டனின் மைடா வேலில் பிறந்தார்.
கோட்பாட்டுக் கணினியியலின் தந்தையாகக் கருதப்படுகிறார். இயற்பியல்
அறிவியல்களிலும் கணிதத்திலும் அவருக்கு இளமையிலேயே ஆர்வம் ஏற்பட்டது.
ஆனால் சிறுவனாக இருந்தபோது படித்த பள்ளிகள் இலக்கியத்திற்கும்
செவ்வியல் படிப்புகளுக்கும் முக்கியத்துவம் கொடுத்ததால், அவருடைய
ஆர்வங்கள் அங்கு நன்கு பிரதிபலிக்கவில்லை. இதன் விளைவாகப் பள்ளியில்
சிறப்பாகச் செயல்படவில்லை; பல ஆசிரியர்களிடமிருந்து கண்டனமும் பெற்றார்.

\olphoto{turing-alan}{ஆலன் டியூரிங்}

% SEGMENT OLP-0631-S02
டியூரிங் கேம்பிரிட்ஜின் கிங்ஸ் கல்லூரியில் இளநிலை மாணவராகக் கணிதம்
படித்தார். கணிக்கத்தக்க சார்பு என்னும் கருத்தைத் துல்லியமாக வரையறுக்கவும்
தீர்மானச் சிக்கல் தீர்மானிக்கவியலாதது என்று நிறுவவும் முயன்று,
1936-இல் இப்போது டியூரிங் பொறி எனப்படும் மாதிரியை உருவாக்கினார்.
ஆனால் அலோன்சோ சர்ச் தமது லாம்டா கலனத்தின் மூலம் இந்த முடிவை
முன்பே நிறுவியிருந்தார். இருப்பினும் சர்ச்சின் முடிவைச் சுட்டியபடியே
டியூரிங்கின் கட்டுரை வெளியானது. சர்ச் அவரைப் பிரின்ஸ்டனுக்கு அழைத்தார்;
1936 முதல் 1938 வரை அங்கிருந்த டியூரிங், சர்ச்சின் மேற்பார்வையில்
முனைவர் பட்டம் பெற்றார்.

% SEGMENT OLP-0631-S03
தருக்கவியலில் ஆர்வம் இருந்தபோதிலும், இயற்பியல் அறிவியல்களில்
டியூரிங்கின் முந்தைய ஆர்வமும் நீடித்தது. இரண்டாம் உலகப் போரின்போது
பிரிட்டனின் பிளெட்ச்லி பார்க்கில் இருந்த மறைகுறியீட்டுப் பகுப்பாய்வுத்
துறையில் பணியாற்றியபோது அவருடைய நடைமுறைத் திறன்கள் பயன்பட்டன.
ஜெர்மன் கடற்படைத் தகவல் தொடர்பில் பயன்படுத்தப்பட்ட எனிக்மா
மறைகுறியீட்டை உடைப்பதில் மையப் பங்கு வகித்தார். புள்ளியியல்,
மறைகுறியீட்டியல் ஆகியவற்றில் டியூரிங்கின் நிபுணத்துவமும் மின்னணு
இயந்திரங்களின் அறிமுகமும் இணைந்து, ``பாம்ப்'' எனப்பட்ட மறைகுறியீட்டை
விலக்கும் இயந்திரத்தை உருவாக்கி, குழுவினர் குறியீட்டை உடைக்க
உதவின. அவருடைய கருத்துகள், பிளெட்ச்லி பார்க்கிலேயே ஜெர்மன்
லோரென்ஸ் மறைகுறியீட்டை உடைக்கப் பயன்படுத்தப்பட்ட, உலகின் முதல்
நிரலாக்கத்தக்க மின்னணுக் கணினி என மூலத்தில் குறிப்பிடப்படும்
கொலோசஸ் உருவாக்கத்திற்கும் உதவின.

% SEGMENT OLP-0631-S04
டியூரிங் ஓர் ஓரினச் சேர்க்கையாளர். இருந்தபோதிலும், 1942-இல்
பிளெட்ச்லி பார்க்கில் தம்முடன் பணியாற்றிய ஜோன் கிளார்க்கைத் திருமணம்
செய்ய முன்மொழிந்தார்; பின்னர் நிச்சயத்தை முறித்துக்கொண்டு தாம்
ஓரினச் சேர்க்கையாளர் என்பதை அவரிடம் தெரிவித்தார். அப்போது
பிரிட்டனில் ஆண்களுக்கு இடையிலான ஓரினச் சேர்க்கைச் செயல்கள்
குற்றங்களாகக் கருதப்பட்டிருந்தாலும், டியூரிங்குக்கு வாழ்நாளில் பல
காதலர்கள் இருந்தனர். 1952-இல் அப்போதைய காதலரின் நண்பர் ஒருவர்
அவருடைய வீட்டில் திருடினார். காவல்துறையில் புகார் செய்தபோது,
ஓரினச் சேர்க்கைச் செயல்களை அரசாங்கம் சட்டப்படி அனுமதிக்கப் போகிறது
என்று தவறாக எண்ணி, தமக்கு ஓரினச் சேர்க்கை உறவு இருப்பதை டியூரிங்
ஒப்புக்கொண்டார். ஆனால் அது உண்மையல்ல; ``கடுமையான ஒழுக்கக்கேடு'' என்ற
குற்றச்சாட்டில் அவர் வழக்கை எதிர்கொண்டார். சிறைக்குச் செல்வதற்குப்
பதிலாகப் பாலுணர்வைக் குறைக்கும் ஹார்மோன் சிகிச்சையைத் தேர்ந்தெடுத்தார்.
1954 ஜூன் 8 அன்று சயனைடு அளவுக்கு அதிகமாக உட்கொண்ட நிலையில்
இறந்திருக்கக் கண்டெடுக்கப்பட்டார்; பெரும்பாலும் அது தற்கொலையாக
இருந்திருக்கலாம். 2013-இல் இரண்டாம் எலிசபெத் அரசி அவருக்கு
அரச மன்னிப்பு வழங்கினார்.

% SEGMENT OLP-0631-S05
\begin{reading}
டியூரிங்கின் விரிவான வாழ்க்கை வரலாற்றுக்கு \citet{Hodges2014}-ஐப்
பார்க்கவும். அவரது வாழ்க்கையும் பணியும் ஊக்கமளித்த \emph{Breaking
  the Code} நாடகம், டியூரிங்காக டெரெக் ஜேக்கபி நடித்த தொலைக்காட்சி
நிகழ்ச்சியாக 1996-இல் தயாரிக்கப்பட்டது. பெனடிக்ட் கம்பர்பேட்ச்,
கீரா நைட்லி நடித்த, அகாடமி விருதுக்குப் பரிந்துரைக்கப்பட்ட
\emph{The Imitation Game} திரைப்படமும் டியூரிங்கின் வாழ்க்கை,
பிளெட்ச்லி பார்க் காலம் ஆகியவற்றைத் தளர்வாக அடிப்படையாகக் கொண்டது
\citep{Imitation2014}.

டியூரிங்கின் வாழ்க்கை, பணி பற்றிய பல பாட்காஸ்ட்கள் \citet{Radiolab2012}
இல் உள்ளன. பிபிசி ஹொரைசனின் \emph{The Strange Life and Death of
  Dr.~Turing} ஆவணப்படத்தை இணையத்தில் பார்க்கலாம் \citep{Sykes1992}.
2012-இல் டியூரிங்கின் நூற்றாண்டை நினைவுகூர உருவாக்கப்பட்ட,
இயங்கும் LEGO டியூரிங் பொறியின் குறும்படம் \citep{Theelen2012}.

டியூரிங் பொறிகளும் தீர்மானச் சிக்கலும் பற்றிய அவருடைய மூலக் கட்டுரை
\citet{Turing1937} ஆகும்.
\end{reading}

\end{document}
